% !TEX program = xelatex
% 컴파일: xelatex squeezing_report_v3.tex  (목차/참조 때문에 2회 실행 권장)
% 한글 출력: kotex + XeLaTeX. Windows 기본 'Malgun Gothic' 사용.
% 그림:
%   squeezing_frontier.png            -- v2 finite-loss ultra frontier
%   squeezing_frontier_od.png         -- v2 in-cell OD diagnostic
%   squeezing_frontier_ideal_qe1.png  -- v3 ideal QE=1, loss=0 wide dense scan
\documentclass[11pt,a4paper]{article}

\usepackage{kotex}
\setmainhangulfont{Malgun Gothic}
\setsanshangulfont{Malgun Gothic}

\usepackage{amsmath,amssymb}
\usepackage{graphicx}
\usepackage{booktabs}
\usepackage{array}
\usepackage{geometry}
\geometry{margin=2.4cm}
\usepackage{caption}
\captionsetup{font=small,labelfont=bf}
\usepackage{enumitem}
\usepackage{xcolor}
\usepackage[hidelinks]{hyperref}
\graphicspath{{./}}
\sloppy

\newcommand{\dB}{\,\mathrm{dB}}
\newcommand{\GHz}{\,\mathrm{GHz}}
\newcommand{\MHz}{\,\mathrm{MHz}}
\newcommand{\degC}{\,^{\circ}\mathrm{C}}

\title{\textbf{$^{85}$Rb D1 이중-$\Lambda$ 사파장혼합(FWM) 트윈빔\\
세기차 스퀴징 최적점: 실험 레퍼런스, in-cell 손실,\\
그리고 이상 검출 한계의 분리}\\[4pt]
\large 보고서 v3}
\author{GABES (Generic Atomic Bloch Equation Solver) 분석 보고서}
\date{2026년 6월 25일}

\begin{document}
\maketitle

\begin{abstract}
\noindent
본 v3 보고서는 v1/v2의 결론을 유지하면서, 새로 수행한 \textbf{이상 검출 한계 스캔}
(\texttt{QE=1}, \texttt{loss=0}, \texttt{fidelity=ultra})과 문헌 레퍼런스 비교를 추가한다.
v2의 finite-loss Ultra 모형에서는 in-cell 전파손실 $\mathrm{OD}(\Delta,T)$가 들어가면서
$\xi(T)$가 저온 이득제한과 고온 손실제한 사이의 내부 최적점을 갖게 되었다. 이 현실적
조건에서는 Sim \textit{et al.}의 $^{85}$Rb D1 double-$\Lambda$ FWM 레퍼런스
($\Delta\approx+0.9\GHz$, $T\approx121\degC$, gain $\approx15$, IDS $\approx-7.8\dB$)와
정성적, 정량적으로 잘 맞는다. 특히 GABES의 Sim detuning readout은 $\xi\approx-7.88\dB$를
주어 실험값과 같은 크기이다.

반면 v3의 이상 검출 스캔은 의도적으로 실험 손실과 검출 floor를 제거한 \emph{source-limit}
질문이다. 넓고 촘촘한 Ultra 격자
$\Delta\in[-6,+8]\GHz$, $T\in[50,220]\degC$에서 전역 최적은
$\Delta=-2.10\GHz$, $T=142.5\degC$, $\delta=-120\MHz$,
$G_s\simeq1.26\times10^4$, $\mathrm{OD}=0$, $\xi=-81.95\dB$로 나타났다.
이 좌표는 문헌의 실험 권고 동작점과 직접 일치하지 않는다. 그러나 이는 모순이 아니다.
검출 손실, pump scatter, seed excess noise, collision noise가 제거되면 목적함수가
"실험에서 가장 강건한 점"이 아니라 "모형 내부에서 거의 무손실이고 균형 이득이 가장 큰 점"으로
바뀌기 때문이다. 따라서 v3의 최종 판정은 다음과 같다:
\textbf{현실 조건의 GABES 최적화는 레퍼런스와 일치하며, \texttt{QE=1} 이상 스캔은 별도의
theoretical source-limit로 분리해서 해석해야 한다.}
\end{abstract}

\tableofcontents
\bigskip

%==================================================================
\section{질문의 분리}
%==================================================================
이 보고서에서 서로 다른 세 질문을 분리한다. 이 분리가 없으면 같은 FWM 모델 안에서도
"최적점"이라는 말이 서로 다른 좌표를 가리킨다.

\begin{enumerate}[nosep]
\item \textbf{실험 레퍼런스 재현}: Sim \textit{et al.}과 Lett 계열 문헌의 85Rb D1
      double-$\Lambda$ FWM squeezing 동작점($T\sim120\degC$, pump detuning
      $\sim+0.8\text{--}1.0\GHz$, IDS $-8\dB$대)을 GABES가 재현하는가?
\item \textbf{현실적 최효율 전선}: 검출효율, optical loss, in-cell 손실을 포함했을 때,
      최적에 거의 도달하면서 온도/이득/OD를 최소화하는 robust operating point는 어디인가?
\item \textbf{이상 source-limit}: \texttt{QE=1}, \texttt{loss=0}으로 검출 floor를 제거하면,
      Ultra 모형 자체가 허용하는 가장 깊은 세기차 스퀴징은 어디인가?
\end{enumerate}

v1은 첫 번째 이상화에서 전역 최대가 검출효율 floor로 collapse하는 문제를 보였다. v2는
in-cell 손실을 넣어 두 번째 질문, 즉 \emph{최효율 전선}을 찾았다. v3는 세 번째 질문을 실제로
실행하고, 그 결과가 문헌의 첫 번째 질문과 왜 달라지는지를 정리한다.

%==================================================================
\section{좌표계와 레퍼런스 설정}
%==================================================================

\subsection{GABES detuning convention}
GABES의 seeded FWM convention은 다음과 같다.
\begin{itemize}[nosep]
\item $\Delta$: 펌프 one-photon detuning. 기준은 $^{85}$Rb D1의
      $F=2\to F'=3$ 선이다.
\item 즉 $\omega_{\mathrm{pump}}=\omega(F=2\to F'=3)+\Delta$.
\item 표준 seeded FWM readout은 $(-)$ Raman branch이며, probe/seed는 대략
      $\Delta-\nu_{\mathrm{HF}}$ 근처에서 스캔된다.
\item 바닥 초미세 분리는 $\nu_{\mathrm{HF}}=3.035732439\GHz$이다.
\end{itemize}

따라서 이번 이상 스캔의 전역 최적 $\Delta=-2.10\GHz$는
$F=2\to F'=3$ 기준으로는 red detuning이지만, $F=3\to F'=3$ 전이 기준으로는
$-2.10+3.0357\simeq+0.936\GHz$ blue detuning이다. 이 숫자는 문헌의
$+0.8\text{--}0.9\GHz$와 닮아 보이지만, 문헌에서 보통 인용되는 Sim 동작점은
$F=2\to F'=3$ 기준 $+0.9\GHz$이므로 같은 물리 좌표라고 섣불리 동일시하면 안 된다.

\subsection{고정 파라미터}
스캔에서 $\Delta,T$를 제외한 seeded FWM 파라미터는 회귀 테스트의
\texttt{CONFIGS['sim\_optimum']}과 동일하게 두었다.

\begin{table}[h]
\centering
\caption{seeded 85Rb D1 FWM 레퍼런스 설정.}
\label{tab:refparams}
\begin{tabular}{lll}
\toprule
항목 & 값 & 비고\\
\midrule
$P_{\mathrm{pump}}$ & $600\,\mathrm{mW}$ & Sim 레퍼런스\\
$P_{\mathrm{seed}}$ & $8\,\mu\mathrm{W}$ & Sim 레퍼런스\\
line-strength residual & $0.74$ & 물리 결합정규화는 엔진 내부\\
셀 길이 $L$ & $12.5\,\mathrm{mm}$ & GABES 상수\\
빔 허리 $w_{\mathrm{pump}}/w_{\mathrm{seed}}$ & $530/330\,\mu\mathrm{m}$ & GABES 상수\\
branch & $-1$ & standard seeded FWM branch\\
\bottomrule
\end{tabular}
\end{table}

현실 조건 스캔은 $\mathrm{QE}=0.92$, post-cell optical loss $5.5\%$를 사용해
$\eta=0.8694$이다. 이상 source-limit 스캔은 런타임에서 $\mathrm{QE}=1$,
post-cell loss $0$으로 두었다. 이때 detection floor는 물리적으로 $-\infty$이며, 코드의
guardrail 표시는 $-3000\dB$이다.

%==================================================================
\section{현실 손실 포함 Ultra 결과: 레퍼런스와 맞는가}
%==================================================================

\subsection{finite-loss Ultra frontier}
v2의 Ultra 스캔은 $\Delta\in[-3.5,+3.0]\GHz$, $T\in[60,150]\degC$에서 수행되었다.
그 결과는 그림~\ref{fig:frontier_v2}와 같다.

\begin{figure}[p]
\centering
\includegraphics[width=\textwidth]{squeezing_frontier.png}
\caption{finite-loss Ultra 스캔(v2). 왼쪽: $\xi(\Delta,T)$ heatmap. 가운데:
최적에 $0.1\dB$ 이내로 도달하는 최저 온도 전선. 오른쪽: 최효율 detuning에서의
$\xi(T)$.}
\label{fig:frontier_v2}
\end{figure}

\begin{table}[h]
\centering
\caption{finite-loss Ultra 스캔의 핵심 좌표.}
\label{tab:v2summary}
\begin{tabular}{lrrrrr}
\toprule
점 & $\Delta$ [GHz] & $T$ [$^\circ$C] & $\xi$ [dB] & $G_s$ & in-cell OD\\
\midrule
전역 최적 & $-2.2$ & $140$ & $-8.8406$ & $9.1\times10^3$ & $0.000$\\
최효율점($0.1\dB$ 이내) & $-1.9$ & $120$ & $-8.7532$ & $32.5$ & $0.0026$\\
Sim detuning readout & $+0.9$ & $120$ & $-7.8766$ & $12.1$ & $0.0336$\\
\bottomrule
\end{tabular}
\end{table}

여기서 중요한 것은 두 가지다. 첫째, 전역 최적은 detection floor
$10\log_{10}(1-\eta)=-8.8406\dB$와 사실상 일치한다. 둘째, 실험적으로 더 의미 있는
최효율점은 $T\approx120\degC$의 중간 이득 영역이다. 이 점은 최적보다 약 $0.09\dB$ 얕지만,
이득과 OD가 훨씬 작아 excess noise에 더 강하다.

\subsection{in-cell 손실이 만드는 내부 최적}
그림~\ref{fig:od_v2}는 v2에서 새로 들어간 in-cell optical depth를 보여준다. 공명 근처에서는
온도가 올라갈수록 gain도 증가하지만 OD도 함께 증가한다. 따라서 $\xi(T)$는 낮은 온도에서
이득 부족, 높은 온도에서 손실 과다라는 두 제한의 경쟁으로 내부 최저점을 갖는다.

\begin{figure}[p]
\centering
\includegraphics[width=\textwidth]{squeezing_frontier_od.png}
\caption{v2의 in-cell OD 진단. 왼쪽: $\mathrm{OD}(\Delta,T)$ map. 오른쪽:
대표 detuning들의 $\xi(T)$ slice. 공명 근처는 고온에서 손실 제한으로 되돌아오며,
투명창 쪽은 floor에 접근한다.}
\label{fig:od_v2}
\end{figure}

이 결과는 문헌과 방향이 같다. 85Rb D1 FWM에서 온도를 올리면 gain은 커지지만, 일정 지점
이후에는 absorption, collision noise, pump scatter, imbalance leakage가 squeezing을 제한한다.
따라서 "온도만 올리면 무한히 좋아진다"는 이전 toy 결론은 Ultra 모형에서는 사라진다.

%==================================================================
\section{이상 검출 한계 스캔: \texttt{QE=1}, \texttt{loss=0}}
%==================================================================

\subsection{실행 조건}
v3에서 새로 실행한 스캔은 실험 권고점을 찾는 목적이 아니라, 검출 손실과 post-cell loss를
완전히 제거했을 때 GABES Ultra 모형의 source-limit 최적점을 찾는 목적이다.

\begin{table}[h]
\centering
\caption{v3 이상 source-limit 스캔 설정.}
\label{tab:idealrun}
\begin{tabular}{ll}
\toprule
항목 & 값\\
\midrule
fidelity & Ultra\\
$\mathrm{QE}$, post-cell loss & $1.0$, $0$\\
$\Delta$ 범위 & $[-6,+8]\GHz$, step $0.05\GHz$\\
$T$ 범위 & $[50,220]\degC$, step $2.5\degC$\\
$\delta$ scan window & $\pm1.0\GHz$\\
coarse points & $201$\\
velocity grid & $2.0\,\mathrm{m/s}$, cutoff $4\sigma$\\
총 격자 & $281\times69=19389$ points\\
실행 시간 & $13012\,\mathrm{s}$\\
\bottomrule
\end{tabular}
\end{table}

고온/고이득 일부 외곽점에서 segment propagation의 overflow warning이 발생했지만,
최종 저장 배열의 $\xi$, $G_s$, $G_c$에는 NaN/Inf가 없었다. 전체 중 $\delta$ window edge에
걸린 점은 $5293$개였으나, 최적점은 edge-limited가 아니었다.

\subsection{전역 최적 결과}
\begin{figure}[p]
\centering
\includegraphics[width=\textwidth]{squeezing_frontier_ideal_qe1.png}
\caption{v3 이상 검출 한계 스캔. 왼쪽: $\xi(\Delta,T)$ heatmap과 작은 원으로 표시한
전역 최적. 가운데: 각 $\Delta$에서 얻을 수 있는 최저 $\xi$와 그때의 온도. 오른쪽:
최적 detuning $\Delta=-2.10\GHz$에서의 $\xi(T)$ slice. y축은 $0$에서 $-100\dB$로 제한했다.}
\label{fig:ideal_qe1}
\end{figure}

\begin{table}[h]
\centering
\caption{v3 이상 source-limit 최적점.}
\label{tab:idealopt}
\begin{tabular}{lr}
\toprule
항목 & 값\\
\midrule
$\Delta$ & $-2.10\GHz$\\
$T$ & $142.5\degC$\\
$\delta$ & $-120\MHz$\\
$\xi$ & $-81.9469\dB$\\
$G_s$ & $12617.07$\\
$G_c$ & $12615.04$\\
in-cell OD & $0.000$\\
edge-limited? & no\\
최적 $0.01\dB$ 이내 점 수 & $1/19389$\\
\bottomrule
\end{tabular}
\end{table}

이 값은 문헌 실험에서 보이는 $-8$에서 $-9\dB$대와 전혀 다르다. 그러나 이는 불일치라기보다
질문의 차이다. \texttt{QE=1}, \texttt{loss=0}에서는 detection floor가 사라지고, 스캔은
기술잡음이 전혀 없는 순수 source model에서 거의 무손실인 투명창과 높은 균형 이득이 만나는
점을 찾는다. 이 경우 $S_{\mathrm{ideal}}\sim 1/G^2$ scaling 때문에 $G\sim10^4$만 되어도
$-80\dB$대가 수학적으로 가능하다.

\subsection{왜 실험 최적점과 같지 않은가}
문헌 실험 최적점은 단순히 가장 큰 intrinsic gain을 찾는 문제가 아니다. 실제 실험에서는
다음 제한들이 동시에 작동한다.

\begin{itemize}[nosep]
\item detection efficiency와 optical loss가 $-8$에서 $-10\dB$대 floor를 만든다.
\item pump scattering과 detector/electronic noise가 low-frequency IDS를 제한한다.
\item seed가 완전한 shot-noise-limited coherent state가 아니면 seed excess noise가 증폭된다.
\item gain이 너무 크면 anti-squeezed sum noise가 imbalance를 통해 difference channel로 샌다.
\item 온도 상승은 collision dephasing, fluorescence, absorption, self-action을 키운다.
\end{itemize}

따라서 이번 $\Delta=-2.10\GHz$, $T=142.5\degC$ 좌표는 "실험 권고점"이 아니라
\textbf{GABES Ultra 모형의 이상 source-limit 좌표}로 라벨링해야 한다. 실험 권고점은 여전히
$T\sim120\degC$ 부근의 finite-loss frontier가 더 적절하다.

%==================================================================
\section{문헌 레퍼런스와의 비교}
%==================================================================

\subsection{85Rb D1 double-$\Lambda$ squeezing 최적점}
직접적으로 비교 가능한 레퍼런스는 다음과 같다.

\begin{table}[h]
\centering
\caption{문헌과 GABES 결과의 비교.}
\label{tab:literature}
\begin{tabular}{p{3.3cm}p{3.4cm}p{3.0cm}p{4.8cm}}
\toprule
출처 & 보고된 조건 & squeezing/gain & GABES와의 관계\\
\midrule
McCormick \textit{et al.} (2007) &
hot Rb vapor FWM &
$-3.5\dB$ measured, $-8.1\dB$ loss-corrected &
초기 실험 기준. GABES finite-loss 결과의 크기($-8\dB$대)와 같은 scale.\\
\addlinespace
Pooser \textit{et al.} (2009) &
85Rb D1, $\Delta\approx+0.8\GHz$ cut, 400 mW pump &
$-8.0(2)\dB$; detection efficiency $\approx0.90$라 $-10\dB$ 상한 &
finite-loss GABES가 $\eta$ floor에 의해 $-8$에서 $-10\dB$대로 제한되는 것과 일치.\\
\addlinespace
Sim \textit{et al.} (2025) &
85Rb D1, $T=121\degC$, $\Delta\approx+0.9\GHz$, $\delta=-8\MHz$ &
gain $\approx15$, IDS $-7.8\dB$, slope estimate $-8.7\dB$ &
GABES Sim detuning readout $T=120\degC$에서 $\xi\approx-7.88\dB$로 잘 맞음.\\
\addlinespace
Allen \textit{et al.} (2026) &
85Rb 12.5 mm cell, $T\sim118\degC$ near optimized detuning &
long-term stable IDS $-7.8\dB$ &
실험적으로 $T\sim118\text{--}121\degC$ 부근이 robust operating point임을 뒷받침.\\
\addlinespace
Jasperse \textit{et al.} (2011) &
85Rb FWM with distributed gain/loss model &
gain과 absorption 경쟁을 analytic model로 설명 &
GABES Ultra의 in-cell loss 채널과 같은 물리 방향.\\
\bottomrule
\end{tabular}
\end{table}

문헌 전체의 메시지는 꽤 일관적이다. 85Rb D1 double-$\Lambda$ FWM에서 좋은 실험 squeezing은
보통 $T\sim120\degC$, pump detuning $+0.8\text{--}1.0\GHz$ 부근, gain $10$에서 $20$ 사이,
측정 IDS $-8\dB$ 전후에서 나온다. 또한 최고 gain tuning이 최고 squeezing tuning과 같지
않고, loss와 absorption이 최적점을 결정한다.

\subsection{일치하는 부분}
GABES finite-loss Ultra 분석은 다음 점에서 문헌과 일치한다.

\begin{enumerate}[nosep]
\item \textbf{온도 scale}: Sim/Allen의 $118\text{--}121\degC$와 GABES 최효율점
      $120\degC$가 일치한다.
\item \textbf{squeezing scale}: Sim의 $-7.8\dB$와 GABES의 Sim detuning readout
      $-7.8766\dB$가 거의 같은 크기다.
\item \textbf{gain scale}: Sim의 gain $\approx15$와 GABES Sim detuning readout의
      $G_s\approx12$가 같은 order이다.
\item \textbf{손실 제한}: Pooser와 Jasperse가 강조한 detection/intrinsic loss 제한이
      GABES의 $\eta$ floor 및 in-cell OD penalty와 같은 방향이다.
\end{enumerate}

\subsection{일치하지 않는 부분과 이유}
v3의 ideal source-limit 최적점은 문헌의 실험 최적점과 일치하지 않는다.

\begin{equation}
\text{문헌 실험 최적: } \Delta\approx+0.8\text{--}+0.9\GHz,\ T\approx120\degC,\ \xi\sim-8\dB
\end{equation}
\begin{equation}
\text{GABES 이상 최적: } \Delta=-2.10\GHz,\ T=142.5\degC,\ \xi=-81.95\dB
\end{equation}

이 차이는 세 가지 이유로 자연스럽다.

\begin{enumerate}[nosep]
\item \textbf{목적함수 차이}: 문헌은 실제 balanced detector에서 관측되는 IDS를 최적화한다.
      v3 ideal scan은 검출 손실과 기술잡음을 제거한 intrinsic source squeezing을 최적화한다.
\item \textbf{좌표계 차이}: $\Delta=-2.10\GHz$는 $F=2\to F'=3$ 기준으로는 red detuning이지만,
      $F=3\to F'=3$ 기준으로는 약 $+0.936\GHz$ blue detuning이다. 숫자만 보고 Sim의
      $+0.9\GHz$와 같은 점으로 해석하면 안 된다.
\item \textbf{calibration domain}: line-strength residual $0.74$는 Sim 레퍼런스 주변에서
      anchor된 값이다. $-2.10\GHz$ 투명창의 초고이득 source-limit를 정량 실험 예측으로 쓰려면
      그 영역의 별도 calibration과 full noise model이 필요하다.
\end{enumerate}

따라서 v3 ideal result는 \emph{모형 내부에서 손실 제한이 제거되면 어디까지 갈 수 있는가}에
대한 답이고, 문헌 검증은 finite-loss Ultra 결과로 해야 한다.

%==================================================================
\section{물리적 해석}
%==================================================================

\subsection{왜 finite-loss에서는 frontier가 중요해지는가}
유한 검출효율에서는
\begin{equation}
S(\eta)=\eta S_{\mathrm{cell}}+(1-\eta)
\end{equation}
이므로 $S_{\mathrm{ideal}}$을 아무리 줄여도 $1-\eta$ 아래로 갈 수 없다. 예를 들어
$\eta=0.8694$이면 floor는 $-8.8406\dB$이다. 이때 $G_s\gtrsim30$이면 이미 floor에
거의 도달하므로, 더 큰 이득은 실험적으로는 이득보다 비용이 크다. 더 큰 $G$는 pump scatter,
anti-squeezed sum noise leakage, thermal/collisional noise를 키운다. 그러므로 실험 권고점은
"가장 깊은 점"이 아니라 "floor에 거의 도달하는 가장 낮은 온도와 작은 OD의 점"이다.

\subsection{왜 ideal scan에서는 $-80\dB$가 나오는가}
loss-free balanced twin beam에서 $G_c\simeq G_s-1$이면
\begin{equation}
S_{\mathrm{ideal}}\approx\frac{5}{8G_s^2}.
\end{equation}
v3 optimum의 $G_s\simeq12617$을 대입하면
\begin{equation}
10\log_{10}\left(\frac{5}{8G_s^2}\right)\approx -84\dB
\end{equation}
scale이 된다. 실제 GABES 값 $-81.95\dB$는 gain imbalance와 full propagation details를
포함한 같은 order의 값이다. 즉 $-81.95\dB$는 숫자상 기괴해 보이지만, \texttt{QE=1},
technical noise $=0$이라는 이상 조건에서는 수식상 이상하지 않다.

\subsection{실험 안티스퀴징은 손실만으로 설명되지 않는다}
손실은 squeezed difference noise를 shot noise 쪽으로 되돌릴 뿐 $0\dB$ 위로 올리지 못한다.
따라서 실험에서 difference channel이 양의 dB로 올라가거나, predicted loss floor보다 훨씬
얕다면 원인은 pure loss가 아니라 excess noise이다. 가장 위험한 항은 큰 이득에서 커지는
anti-squeezed sum noise가 imbalance를 통해 difference readout으로 새는 경로다. 이 때문에
문헌의 $T\sim120\degC$ 동작점은 단순히 "조금 덜 깊은 점"이 아니라 실제 장치에서는
가장 안정한 선택이다.

%==================================================================
\section{결론}
%==================================================================

\begin{enumerate}[nosep]
\item \textbf{레퍼런스 일치 여부}: finite-loss Ultra GABES 결과는 85Rb D1 double-$\Lambda$
      FWM squeezing 문헌과 잘 맞는다. Sim detuning readout의 $\xi\approx-7.88\dB$는
      Sim \textit{et al.}의 $-7.8\dB$와 직접 비교 가능한 수준이다.
\item \textbf{in-cell 손실의 역할}: v2 이후 모형에서는 온도를 올린다고 스퀴징이 무한정
      좋아지지 않는다. OD가 커지는 detuning에서는 $\xi(T)$가 내부 최저점을 갖고 이후
      악화된다.
\item \textbf{실험 권고점}: 현실 조건에서는 $\Delta\approx-1.9\GHz$, $T\approx120\degC$의
      finite-loss frontier가 가장 실용적이다. Sim의 $+0.9\GHz$ 동작점도 같은 dB scale에서
      레퍼런스를 잘 재현한다.
\item \textbf{이상 source-limit}: \texttt{QE=1}, \texttt{loss=0} wide-dense Ultra 스캔의 전역
      최적은 $\Delta=-2.10\GHz$, $T=142.5\degC$, $\xi=-81.95\dB$이다. 이는 실험 권고점이
      아니라 model-internal theoretical source-limit이다.
\item \textbf{해석 원칙}: 앞으로 보고서나 UI에 두 값을 함께 제시할 때는
      \emph{reference-calibrated experimental operating point}와
      \emph{ideal source-limit optimum}을 명확히 분리해야 한다.
\end{enumerate}

\noindent\textbf{한 줄 요약.} \emph{GABES는 문헌의 85Rb D1 FWM squeezing 동작점을 현실 손실
조건에서 잘 재현한다. 이번 $-81.95\dB$ 결과는 문헌과 충돌하는 실험 예측이 아니라,
손실과 검출 floor를 제거했을 때 Ultra 모형이 보여주는 별도의 이상 source-limit이다.}

\appendix
%==================================================================
\section{수치 재현 정보}
%==================================================================

\begin{itemize}[nosep]
\item finite-loss v2 산출물:
      \texttt{analysis/squeezing/squeezing\_frontier.npz},
      \texttt{docs/squeezing\_report/squeezing\_frontier.png},
      \texttt{docs/squeezing\_report/squeezing\_frontier\_od.png}.
\item ideal v3 산출물:
      \texttt{analysis/squeezing/squeezing\_frontier\_ultra\_ideal\_qe1\_wide\_dense/}
      \texttt{squeezing\_frontier.npz}.
\item ideal v3 그림:
      \texttt{docs/squeezing\_report/squeezing\_frontier\_ideal\_qe1.png}.
\item ideal v3 finite check: $\xi$, $G_s$, $G_c$ NaN/Inf $0$개. edge-limited point
      $5293/19389$, optimum은 edge-limited 아님.
\item ideal v3 warning: 고온/고이득 일부 외곽점에서 matrix propagation overflow warning이
      발생했으나, 최종 저장 배열은 finite이고 optimum은 warning 외곽영역이 아니다.
\end{itemize}

%==================================================================
\section{참고문헌}
%==================================================================

\begin{enumerate}[label={[\arabic*]},leftmargin=*]
\item G. Sim, H. Kim, H. S. Moon, ``Intensity-difference squeezing from four-wave mixing
      in hot $^{85}$Rb and $^{87}$Rb atoms in single diode laser pumping system,''
      \textit{Scientific Reports} \textbf{15}, 7727 (2025).
      \url{https://www.nature.com/articles/s41598-025-86479-w}
\item C. F. McCormick, V. Boyer, E. Arimondo, P. D. Lett,
      ``Strong relative intensity squeezing by four-wave mixing in rubidium vapor,''
      \textit{Optics Letters} \textbf{32}, 178--180 (2007).
      \url{https://doi.org/10.1364/OL.32.000178}
\item R. C. Pooser, A. M. Marino, V. Boyer, K. M. Jones, P. D. Lett,
      ``Quantum correlated light beams from non-degenerate four-wave mixing in an atomic vapor:
      the D1 and D2 lines of $^{85}$Rb and $^{87}$Rb,''
      \textit{Optics Express} \textbf{17}, 16722--16730 (2009).
      \url{https://doi.org/10.1364/OE.17.016722}
\item M. Jasperse, L. D. Turner, R. E. Scholten,
      ``Relative intensity squeezing by four-wave mixing with loss: an analytic model and
      experimental diagnostic,'' \textit{Optics Express} \textbf{19}, 3765--3774 (2011).
      \url{https://arxiv.org/abs/1012.3482}
\item C. H. Allen, M. San Soucie, R. J. Benson, O.-M. Glezakou-Elbert, R. Jimenez,
      J. Morrell-Falvey, Y.-Z. Ma,
      ``Long-term stabilization of intensity-difference squeezing from four-wave mixing in
      rubidium vapor,'' \textit{Optics Express} \textbf{34}, 19382--19399 (2026).
      \url{https://doi.org/10.1364/OE.600911}
\end{enumerate}

\end{document}
